\section{结果与分析}

\subsection{模型性能比较}

表 \ref{tab:model_comparison} 汇总了 Logistic Regression 基准模型
以及 7 个候选模型在验证集和测试集上的性能。
所有模型使用相同的预处理结果和固定数据划分，
并仅使用训练集进行拟合，
表中结果可用于比较不同模型在基础参数设置下的表现。

\begin{table*}[t]
  \caption{模型性能比较表}
  \label{tab:model_comparison}
  \centering
  \scriptsize
  \begin{tabular}{lcccccccccc}
    \toprule
    模型
    & Val Acc. & Val Prec. & Val Recall & Val F1 & Val AUC
    & Test Acc. & Test Prec. & Test Recall & Test F1 & Test AUC \\
    \midrule
    Logistic Regression
    & 0.7642 & 0.6346 & 0.7674 & 0.6947 & 0.8709
    & 0.7273 & 0.5938 & 0.7037 & 0.6441 & 0.8154 \\
    KNN
    & 0.7805 & 0.7105 & 0.6279 & 0.6667 & 0.8581
    & 0.7273 & 0.6154 & 0.5926 & 0.6038 & 0.7675 \\
    Decision Tree
    & 0.7886 & 0.7429 & 0.6047 & 0.6667 & 0.7461
    & 0.6948 & 0.5714 & 0.5185 & 0.5437 & 0.6543 \\
    Random Forest
    & 0.7967 & 0.7368 & 0.6512 & 0.6914 & 0.8365
    & 0.7338 & 0.6444 & 0.5370 & 0.5859 & 0.8275 \\
    SVM
    & 0.7967 & 0.6667 & 0.8372 & 0.7423 & 0.8759
    & 0.7468 & 0.6087 & 0.7778 & 0.6829 & 0.8157 \\
    Gradient Boosting
    & 0.7236 & 0.6098 & 0.5814 & 0.5952 & 0.8311
    & 0.7013 & 0.5952 & 0.4630 & 0.5208 & 0.7959 \\
    XGBoost
    & 0.7561 & 0.6857 & 0.5581 & 0.6154 & 0.7985
    & 0.7468 & 0.6596 & 0.5741 & 0.6139 & 0.7967 \\
    LightGBM
    & 0.7724 & 0.6923 & 0.6279 & 0.6585 & 0.8308
    & 0.7273 & 0.6304 & 0.5370 & 0.5800 & 0.8002 \\
    \bottomrule
  \end{tabular}
\end{table*}

\begin{figure*}[t]
  \centering
  \includegraphics[width=0.95\textwidth]{model_comparison.png}
  \caption{不同模型在测试集上的 Accuracy、Recall、F1-score 和 ROC AUC 对比}
  \label{fig:model_comparison}
\end{figure*}

图 \ref{fig:model_comparison} 进一步展示了各模型在测试集上的
Accuracy、Recall、F1-score 和 ROC AUC。
由于糖尿病风险预测不能只关注总体预测正确率，
本文在比较候选模型时重点考察 Recall 和 F1-score，
并使用 ROC AUC 衡量模型对正负样本的整体排序能力。

在基础参数设置下，SVM 与 XGBoost 的测试集 Accuracy 均为 0.7468，
并列所有模型中的最高值。
但 SVM 的 Recall 为 0.7778，F1-score 为 0.6829，
二者均为所有候选模型中的最高值。
相比 Logistic Regression 基准模型，
SVM 的 Recall 从 0.7037 提升至 0.7778，
F1-score 从 0.6441 提升至 0.6829。
这说明 SVM 能够识别更多阳性样本，
在筛查场景下表现出较好的应用潜力。

Random Forest 的测试集 ROC AUC 为 0.8275，
是所有模型中的最高值，
说明其概率排序能力较好。
但其 Recall 仅为 0.5370，
低于 SVM 和 Logistic Regression，
表明默认分类阈值下仍可能漏判较多阳性样本。
该结果说明，Random Forest 具有较好的整体区分能力，
但仍需要结合参数优化和阈值调整改善阳性样本识别效果。

XGBoost 的测试集 Accuracy 为 0.7468，
Precision 为 0.6596，
但 Recall 仅为 0.5741，
说明其预测相对保守。
LightGBM 的各项指标较为均衡，
但没有超过 SVM 的 Recall 和 F1-score，
也没有超过 Random Forest 的 ROC AUC。
Decision Tree 和 Gradient Boosting 在默认参数下表现相对较弱。
综合考虑筛查任务中漏判阳性样本的风险，
SVM 和 Random Forest 可作为后续优化阶段的重点候选模型，
并应结合 Recall、F1-score、ROC AUC 和分类阈值进行综合选择。

\subsection{误判与 ROC 分析}

当前基准模型在验证集上的混淆矩阵为：
TN=61、FP=19、FN=10、TP=33；
在测试集上为：
TN=74、FP=26、FN=16、TP=38。
从筛查任务角度看，16 个测试集假阴性意味着真实阳性样本被判断为阴性，
这是后续模型优化时应重点减少的错误类型。

\placeholder{补充最佳模型的混淆矩阵热力图和 ROC 曲线。
如果已生成多模型 ROC 曲线，可在此比较各模型 AUC 和曲线形态。}

\subsection{特征解释}

当前可用的解释结果来自 Logistic Regression 基准模型。
图 \ref{fig:baseline_feature_importance} 展示了标准化系数绝对值排序。
Glucose、BMI、DiabetesPedigreeFunction、Pregnancies 和 Age 位于前五位；
其中 Glucose、BMI 和 Age 与相关研究中强调的重要因素基本一致
\cite{LiJiaSiJiYuJiQiXueXiDeTangNiaoBingYuCeJiSHAPTeZhengFenXi2023,
ChengLuPingJiYuJiQiXueXiDe2XingTangNiaoBingFaBingYinSuFenXiYuYuCe2025}。
需要注意的是，线性模型系数只反映当前数据划分和模型假设下的统计关联，
不能解释为医学因果关系。

\begin{figure}[t]
  \centering
  \includegraphics[width=\linewidth]{baseline_feature_importance.png}
  \caption{Logistic Regression 基准模型特征影响排序}
  \label{fig:baseline_feature_importance}
\end{figure}

\placeholder{若最终模型为 Random Forest、XGBoost 或其他非线性模型，
在此替换或补充特征重要性、置换重要性或 SHAP 图。}

\subsection{系统演示}

为展示模型在单样本预测场景中的应用效果，
本文进一步构建了基于 Gradio 的糖尿病风险预测问诊系统。
系统输入与 Pima 数据集的 8 个特征保持一致，
包括 Pregnancies、Glucose、BloodPressure、SkinThickness、
Insulin、BMI、DiabetesPedigreeFunction 和 Age。
在交互界面中，这些变量被转换为易于理解的中文表单项，
用户可以根据个体情况填写或调整相应数值。
系统同时保留典型样例入口，
用于在展示时快速说明低风险和高风险输入在模型输出上的差异。

系统预测流程严格复用实验阶段的预处理协议。
对于网页端提交的原始指标，系统首先按照训练阶段的规则处理不合理 0 值，
再使用训练集拟合得到的中位数填充器、IQR 裁剪边界和标准化器进行变换。
完成预处理后，样本被送入实验阶段选定的分类模型，
模型输出糖尿病阳性的预测概率。
由于预处理对象和模型对象均由训练阶段保存得到，
系统端推理与离线实验保持了相同的特征顺序和数据变换方式，
从而提高了系统展示结果与实验结果之间的一致性。

系统输出由三部分组成。
第一部分是阳性预测概率，用于表示模型认为该输入样本属于糖尿病阳性类别的可能性；
第二部分是风险等级，
本文按概率阈值将低于 0.35 的样本标记为低风险，
将位于 $[0.35,0.65)$ 的样本标记为中风险，
将不低于 0.65 的样本标记为高风险；
第三部分是简要解释文本，
用于提示血糖、BMI、年龄、血压或糖尿病家族遗传指数等较突出的输入因素。
因此，系统展示并不只给出二分类标签，
而是同时呈现概率、等级和输入因素说明，
使预测结果更适合课程展示和非技术用户理解。

该系统仅用于机器学习课程项目展示，
其输出不能替代医生诊断或真实临床决策。
